\chapter{2015年浙江卷（文）}

\section{单选题}

\begin{problem}
已知集合 \(P=\{x\mid x^2-2x\geqslant3\}\)，\(Q=\{x\mid 2<x<4\}\)，则 \(P\cap Q=\)（\quad）

\choices
    {\([3,4)\)}
    {\((2,3]\)}
    {\((-1,2)\)}
    {\((-1,3]\)}
\end{problem}

\begin{answer}
A
\end{answer}

\begin{solution}
由 \(x^2-2x\geqslant3\)，得 \(x^2-2x-3\geqslant0\)，即 \((x-3)(x+1)\geqslant0\)，所以 \(P=(-\infty,-1]\cup[3,+\infty)\). 与 \(Q=(2,4)\) 取交集，得 \(P\cap Q=[3,4)\)，故选 A.
\end{solution}

\begin{problem}
某几何体的三视图如图所示（单位：\(\mathrm{cm}\)），则该几何体的体积是（\quad）

\bitmapfigure[max width=0.60\linewidth]{img/2015/zhejiang_liberal/q2_fig1.png}

\choices
    {\(8\mathrm{~cm}^3\)}
    {\(12\mathrm{~cm}^3\)}
    {\(\frac{32}{3}\mathrm{~cm}^3\)}
    {\(\frac{40}{3}\mathrm{~cm}^3\)}
\end{problem}

\begin{answer}
C
\end{answer}

\begin{solution}
由三视图可知，该几何体由一个棱长为 \(2\mathrm{~cm}\) 的正方体和一个底面边长为 \(2\mathrm{~cm}\)、高为 \(2\mathrm{~cm}\) 的正四棱锥组成. 因此体积为
\[
V=2^3+\frac{1}{3}\times2^2\times2=8+\frac{8}{3}=\frac{32}{3}\mathrm{~cm}^3
\]
故选 C.
\end{solution}

\begin{problem}
设 \(a\)，\(b\) 是实数，则“\(a+b>0\)”是“\(ab>0\)”的（\quad）

\choices
    {充分不必要条件}
    {必要不充分条件}
    {充分必要条件}
    {既不充分也不必要条件}
\end{problem}

\begin{answer}
D
\end{answer}

\begin{solution}
分别检验两个方向：取 \(a=3\)，\(b=-2\)，有 \(a+b=1>0\)，而 \(ab=-6<0\)，所以 \(a+b>0\) 不是 \(ab>0\) 的充分条件；取 \(a=-1\)，\(b=-2\)，有 \(ab=2>0\)，而 \(a+b=-3<0\)，所以 \(a+b>0\) 不是 \(ab>0\) 的必要条件. 因此两者既不充分也不必要，故选 D.
\end{solution}

\begin{problem}
设 \(\alpha\)，\(\beta\) 是两个不同的平面，\(l\)，\(m\) 是两条不同的直线，且 \(l\subset\alpha\)，\(m\subset\beta\)（\quad）

\choices
    {若 \(l\perp\beta\)，则 \(\alpha\perp\beta\)}
    {若 \(\alpha\perp\beta\)，则 \(l\perp m\)}
    {若 \(l\parallel\beta\)，则 \(\alpha\parallel\beta\)}
    {若 \(\alpha\parallel\beta\)，则 \(l\parallel m\)}
\end{problem}

\begin{answer}
A
\end{answer}

\begin{solution}
若直线 \(l\perp\beta\)，且 \(l\subset\alpha\)，则平面 \(\alpha\) 含有一条垂直于平面 \(\beta\) 的直线，由平面与平面垂直的判定定理，得 \(\alpha\perp\beta\)，所以选项 A 正确.

其余选项均不成立：若 \(\alpha\perp\beta\)，可取 \(l\)，\(m\) 均平行于两平面的交线，此时不垂直；若 \(l\parallel\beta\)，\(\alpha\) 仍可能与 \(\beta\) 相交，不能推出 \(\alpha\parallel\beta\)；若 \(\alpha\parallel\beta\)，分别位于两个平行平面内的直线可能异面，并不一定平行. 故选 A.
\end{solution}

\begin{problem}
函数 \(f(x)=\left(x-\frac{1}{x}\right)\cos x\left(-\pi\leqslant x\leqslant \pi\text{且}\,x\neq0\right)\) 的图象可能为（\quad）

\choices
    {\choicebitmap[width=0.15\paperwidth]{img/2015/zhejiang_liberal/q5_opt_a.png}}
    {\choicebitmap[width=0.15\paperwidth]{img/2015/zhejiang_liberal/q5_opt_b.png}}
    {\choicebitmap[width=0.15\paperwidth]{img/2015/zhejiang_liberal/q5_opt_c.png}}
    {\choicebitmap[width=0.15\paperwidth]{img/2015/zhejiang_liberal/q5_opt_d.png}}
\end{problem}

\begin{answer}
D
\end{answer}

\begin{solution}
因为 \(x-\dfrac{1}{x}\) 是奇函数，\(\cos x\) 是偶函数，所以 \(f(x)\) 是奇函数，图象关于原点对称，排除选项 A、B. 又当 \(x\to0^+\) 时，\(x-\dfrac{1}{x}\to-\infty\)，且 \(\cos x\to1\)，故 \(f(x)\to-\infty\)；当 \(x\to0^-\) 时，\(f(x)\to+\infty\). 在 \(0<x<1\) 时，\(x-1/x<0\) 且 \(\cos x>0\)，故 \(f(x)<0\)；在 \(1<x<\pi/2\) 时，\(f(x)>0\)；在 \(\pi/2<x\leqslant\pi\) 时，\(f(x)<0\)，且 \(f(\pi)=1/\pi-\pi<0\). 右半支依次经过负值、正值、负值区域，符合选项 D，故选 D.
\end{solution}

\begin{problem}
有三个房间需要粉刷，粉刷方案要求：每个房间只用一种颜色，且三个房间颜色各不相同. 已知三个房间的粉刷面积（单位：\(\mathrm{m}^2\)）分别为 \(x\)，\(y\)，\(z\)，且 \(x<y<z\)，三种颜色涂料的粉刷费用（单位：元/\(\mathrm{m}^2\)）分别为 \(a\)，\(b\)，\(c\)，且 \(a<b<c\). 在不同的方案中，最低的总费用（单位：元）是（\quad）

\choices
    {\(ax+by+cz\)}
    {\(az+by+cx\)}
    {\(ay+bz+cx\)}
    {\(ay+bx+cz\)}
\end{problem}

\begin{answer}
B
\end{answer}

\begin{solution}
比较任意两个相邻的分配. 若面积 \(u<v\)，费用 \(r<s\)，把低费用 \(r\) 分配给较大面积 \(v\)、高费用 \(s\) 分配给较小面积 \(u\) 时，总费用为 \(vr+us\)；交换这两个费用后的总费用为 \(ur+vs\)，且
\[
(ur+vs)-(vr+us)=(v-u)(s-r)>0
\]
所以低费用应分配给较大面积. 因此总费用最小时，应把最低单价 \(a\) 分给最大面积 \(z\)，把最高单价 \(c\) 分给最小面积 \(x\)，中间单价 \(b\) 分给 \(y\). 最低总费用为 \(az+by+cx\)，故选 B.
\end{solution}

\begin{problem}
如图，斜线段 \(AB\) 与平面 \(\alpha\) 所成的角为 \(60^\circ\)，\(B\) 为斜足，平面 \(\alpha\) 上的动点 \(P\) 满足 \(\angle PAB=30^\circ\)，则点 \(P\) 的轨迹是（\quad）

\bitmapfigure[max width=0.52\linewidth]{img/2015/zhejiang_liberal/q7_fig1.png}

\choices
    {直线}
    {抛物线}
    {椭圆}
    {双曲线的一支}
\end{problem}

\begin{answer}
C
\end{answer}

\begin{solution}
从点 \(A\) 出发且与 \(AB\) 成 \(30^\circ\) 的所有射线构成以 \(AB\) 为轴、半顶角为 \(30^\circ\) 的圆锥面. 平面 \(\alpha\) 与圆锥轴 \(AB\) 的夹角为 \(60^\circ\)，大于圆锥的半顶角而小于直角，因此该平面截圆锥得到椭圆. 故选 C.
\end{solution}

\begin{problem}
设实数 \(a\)，\(b\)，\(t\) 满足 \(|a+1|=|\sin b|=t\)，则下列说法中正确的是（\quad）

\choices
    {若 \(t\) 确定，则 \(b^2\) 唯一确定}
    {若 \(t\) 确定，则 \(a^2+2a\) 唯一确定}
    {若 \(t\) 确定，则 \(\sin\dfrac{b}{2}\) 唯一确定}
    {若 \(t\) 确定，则 \(a^2+a\) 唯一确定}
\end{problem}

\begin{answer}
B
\end{answer}

\begin{solution}
由 \(|a+1|=t\)，平方得 \((a+1)^2=t^2\)，从而 \(a^2+2a=t^2-1\). 因此当 \(t\) 确定时，\(a^2+2a\) 唯一确定.

而由 \(|\sin b|=t\) 只能确定 \(\sin^2b=t^2\). 例如取 \(t=1/2\)，\(b=\pi/6\) 或 \(b=5\pi/6\)，两者均满足条件但 \(b^2\) 和 \(\sin(b/2)\) 不同；同样，取 \(t=1\)，则 \(a=0\) 或 \(a=-2\)，而 \(a^2+a\) 分别为 \(0\) 和 \(2\). 故选 B.
\end{solution}

\section{填空题}

\begin{problem}
计算 \(\log_2\dfrac{\sqrt{2}}{2}=\)\fillinblank{}，\(2^{\log_2 3+\log_4 3}=\)\fillinblank{}.
\end{problem}

\begin{answer}
\(-\frac{1}{2}\)，\(3\sqrt{3}\)
\end{answer}

\begin{solution}
因为 \(\dfrac{\sqrt{2}}{2}=2^{-1/2}\)，所以 \(\log_2\dfrac{\sqrt{2}}{2}=-\dfrac12\). 又 \(\log_4 3=\dfrac12\log_2 3\)，故
\[
2^{\log_2 3+\log_4 3}=2^{\frac32\log_2 3}=\left(2^{\log_2 3}\right)^{3/2}=3\sqrt{3}
\]
\end{solution}

\begin{problem}
已知 \(\{a_n\}\) 是等差数列，公差 \(d\) 不为零. 若 \(a_2\)，\(a_3\)，\(a_7\) 成等比数列，且 \(2a_1+a_2=1\)，则 \(a_1=\)\fillinblank{}，\(d=\)\fillinblank{}.
\end{problem}

\begin{answer}
\(\frac{2}{3}\)，\(-1\)
\end{answer}

\begin{solution}
设 \(a_n=a_1+(n-1)d\). 由 \(a_2\)，\(a_3\)，\(a_7\) 成等比数列，得
\[
(a_1+2d)^2=(a_1+d)(a_1+6d)
\]
整理为 \(d(3a_1+2d)=0\). 因 \(d\neq0\)，所以 \(3a_1+2d=0\). 另一方面，\(2a_1+a_2=3a_1+d=1\). 联立得 \(a_1=\dfrac23\)，\(d=-1\).
\end{solution}

\begin{problem}
函数 \(f(x)=\sin^2x+\sin x\cos x+1\) 的最小正周期是\fillinblank{}，最小值是\fillinblank{}.
\end{problem}

\begin{answer}
\(\pi\)，\(\frac{3-\sqrt{2}}{2}\)
\end{answer}

\begin{solution}
利用倍角公式，得
\[
f(x)=\frac{1-\cos2x}{2}+\frac{\sin2x}{2}+1=\frac32+\frac{\sin2x-\cos2x}{2}=\frac32+\frac{\sqrt2}{2}\sin\left(2x-\frac\pi4\right)
\]
因此最小正周期为 \(\pi\)，振幅为 \(\dfrac{\sqrt2}{2}\)，最小值为 \(\dfrac32-\dfrac{\sqrt2}{2}=\dfrac{3-\sqrt2}{2}\).
\end{solution}

\begin{problem}
已知函数 \(f(x)=\begin{cases}
x^2,&x\leqslant1,\\
x+\dfrac{6}{x}-6,&x>1.
\end{cases}\) 则 \(f(f(-2))=\)\fillinblank{}，\(f(x)\) 的最小值是\fillinblank{}.
\end{problem}

\begin{answer}
\(-\frac{1}{2}\)，\(2\sqrt{6}-6\)
\end{answer}

\begin{solution}
因为 \(-2\leqslant1\)，所以 \(f(-2)=(-2)^2=4\). 再因 \(4>1\)，有
\[
f(f(-2))=f(4)=4+\frac64-6=-\frac12
\]

当 \(x\leqslant1\) 时，\(f(x)=x^2\geqslant0\). 当 \(x>1\) 时，由基本不等式
\[
x+\frac6x\geqslant2\sqrt6
\]
且等号在 \(x=\sqrt6>1\) 时成立，所以这一段的最小值为 \(2\sqrt6-6<0\). 故全函数的最小值为 \(2\sqrt6-6\).
\end{solution}

\begin{problem}
已知 \(\bs{e}_1\)，\(\bs{e}_2\) 是平面单位向量，且 \(\bs{e}_1\cdot\bs{e}_2=\frac12\). 若平面向量 \(\bs{b}\) 满足 \(\bs{b}\cdot\bs{e}_1=\bs{b}\cdot\bs{e}_2=1\)，则 \(|\bs{b}|=\)\fillinblank{}.
\end{problem}

\begin{answer}
\(\frac{2\sqrt{3}}{3}\)
\end{answer}

\begin{solution}
由 \(\bs{b}\cdot\bs{e}_1=\bs{b}\cdot\bs{e}_2\)，得 \(\bs{b}\cdot(\bs{e}_1-\bs{e}_2)=0\). 又 \(\bs{e}_1-\bs{e}_2\) 与 \(\bs{e}_1+\bs{e}_2\) 的数量积为 \(|\bs{e}_1|^2-|\bs{e}_2|^2=0\)，且它们均为非零向量，所以 \(\bs{b}\) 与 \(\bs{e}_1-\bs{e}_2\) 垂直时，\(\bs{b}\) 可设为 \(\lambda(\bs{e}_1+\bs{e}_2)\). 代入 \(\bs{b}\cdot\bs{e}_1=1\)，得 \(\lambda(1+\frac12)=1\)，即 \(\lambda=\frac23\). 因此
\[
|\bs{b}|=\frac23|\bs{e}_1+\bs{e}_2|=\frac23\sqrt{|\bs{e}_1|^2+|\bs{e}_2|^2+2\bs{e}_1\cdot\bs{e}_2}=\frac23\sqrt3=\frac{2\sqrt3}{3}
\]
\end{solution}

\begin{problem}
已知实数 \(x\)，\(y\) 满足 \(x^2+y^2\leqslant1\)，则 \(|2x+y-4|+|6-x-3y|\) 的最大值是\fillinblank{}.
\end{problem}

\begin{answer}
15
\end{answer}

\begin{solution}
由柯西不等式，\(2x+y\leqslant\sqrt5\)，所以 \(2x+y-4<0\)；同理 \(x+3y\leqslant\sqrt{10}<6\)，所以 \(6-x-3y>0\). 原式化为
\[
(4-2x-y)+(6-x-3y)=10-3x-4y
\]
再由柯西不等式，\(-3x-4y\leqslant\sqrt{3^2+4^2}\sqrt{x^2+y^2}\leqslant5\)，故原式不超过 \(15\). 当 \((x,y)=(-\frac35,-\frac45)\) 时满足 \(x^2+y^2=1\)，且 \(-3x-4y=5\)，取得等号. 因此最大值为 \(15\).
\end{solution}

\begin{problem}
椭圆 \(\frac{x^2}{a^2}+\frac{y^2}{b^2}=1\)（\(a>b>0\)）的右焦点 \(F(c,0)\) 关于直线 \(y=\frac{b}{c}x\) 的对称点 \(Q\) 在椭圆上，则椭圆的离心率是\fillinblank{}.
\end{problem}

\begin{answer}
\(\frac{\sqrt{2}}{2}\)
\end{answer}

\begin{solution}
因 \(a>b>0\)，有 \(c^2=a^2-b^2>0\). 令 \(r=\dfrac bc>0\)，则 \(a^2=c^2+b^2=c^2(1+r^2)\). 直线 \(y=\dfrac bc x\) 的方向向量为 \((1,r)\)，把点 \(F(c,0)\) 关于该直线反射，得到
\[
Q\left(c\frac{1-r^2}{1+r^2},\,c\frac{2r}{1+r^2}\right)
\]
将其代入椭圆方程，得
\[
\frac{(1-r^2)^2}{(1+r^2)^3}+\frac{4}{(1+r^2)^2}=1
\]
令 \(s=r^2>0\)，化简为 \(s^3+2s^2+s-4=0\)，即
\[
(s-1)(s^2+3s+4)=0
\]
因 \(s^2+3s+4>0\)，故 \(s=1\)，即 \(b=c\). 于是 \(a^2=b^2+c^2=2c^2\)，所以离心率
\[
e=\frac ca=\frac{\sqrt2}{2}
\]
\end{solution}

\section{解答题}

\begin{problem}
在 \(\triangle ABC\) 中，内角 \(A\)，\(B\)，\(C\) 所对的边分别为 \(a\)，\(b\)，\(c\). 已知 \(\tan\left(\frac\pi4+A\right)=2\).

\begin{enumerate}
\item 求 \(\frac{\sin 2A}{\sin 2A+\cos^2 A}\) 的值；
\item 若 \(B=\frac\pi4\)，\(a=3\)，求 \(\triangle ABC\) 的面积.
\end{enumerate}
\end{problem}

\begin{solution}
\begin{enumerate}
\item 因为 \(0<A<\pi\)，且 \(\tan\left(\frac\pi4+A\right)=2\)，由正切的和角公式得
\[
\frac{1+\tan A}{1-\tan A}=2
\]
因此 \(\tan A=\frac13\). 于是
\(\sin2A=\frac{2\tan A}{1+\tan^2A}=\frac35\)，\(\cos^2A=\frac{1}{1+\tan^2A}=\frac9{10}\)
从而
\[
\frac{\sin2A}{\sin2A+\cos^2A}=\frac{\frac35}{\frac35+\frac9{10}}=\frac25
\]

\item 由 \(\tan A=\frac13\)，得 \(\sin A=\frac1{\sqrt{10}}\)，\(\cos A=\frac3{\sqrt{10}}\). 由正弦定理，
\[
b=\frac{a\sin B}{\sin A}=\frac{3\cdot\frac{\sqrt2}{2}}{\frac1{\sqrt{10}}}=3\sqrt5
\]
又 \(C=\pi-A-B\)，所以 \(\sin C=\sin(A+B)=\sin A\cos B+\cos A\sin B=\frac2{\sqrt5}\). 因此
\[
S_{\triangle ABC}=\frac12ab\sin C=\frac12\times3\times3\sqrt5\times\frac2{\sqrt5}=9
\]
\end{enumerate}
\end{solution}

\begin{problem}
已知数列 \(\{a_n\}\) 和 \(\{b_n\}\) 满足 \(a_1=2\)，\(b_1=1\)，\(a_{n+1}=2a_n\)（\(n\in\N^{*}\)），\(b_1+\frac{1}{2}b_2+\frac{1}{3}b_3+\dots+\frac{1}{n}b_n=b_{n+1}-1\)（\(n\in\N^{*}\)）.

\begin{enumerate}
\item 求 \(a_n\) 与 \(b_n\)；
\item 记数列 \(\{a_nb_n\}\) 的前 \(n\) 项和为 \(T_n\)，求 \(T_n\).
\end{enumerate}
\end{problem}

\begin{solution}
\begin{enumerate}
\item 由 \(a_1=2\) 及 \(a_{n+1}=2a_n\)，知 \(\{a_n\}\) 是首项为 \(2\)、公比为 \(2\) 的等比数列，所以 \(a_n=2^n\).

令题给的关于 \(b_n\) 的等式分别取 \(n\) 和 \(n-1\)，两式相减得
\[
\frac1n b_n=b_{n+1}-b_n
\]
当 \(n=1\) 时，原等式给出 \(b_2-1=b_1=1\)，所以 \(b_2=2\). 对 \(n\geqslant2\)，上式可写为 \(b_{n+1}=\dfrac{n+1}{n}b_n\). 因而由数学归纳法，\(b_n=n\).

\item 由上问，\(a_nb_n=n2^n\)，故
\[
T_n=\sum_{k=1}^n k2^k
\]
将上式乘以 \(2\)，并与原式相减，得
\[
T_n-2T_n=\sum_{k=1}^n k2^k-\sum_{k=1}^n k2^{k+1}=2+2^2+2^3+\cdots+2^n-n2^{n+1}
=2^{n+1}-2-n2^{n+1}
\]
所以 \(-T_n=(1-n)2^{n+1}-2\)，即 \(T_n=(n-1)2^{n+1}+2\).
\end{enumerate}
\end{solution}

\begin{problem}
如图，在三棱柱 \(ABC-A_1B_1C_1\) 中，\(\angle BAC=90^\circ\)，\(AB=AC=2\)，\(A_1A=4\)，\(A_1\) 在底面 \(ABC\) 的射影为 \(BC\) 的中点，\(D\) 是 \(B_1C_1\) 的中点.

\begin{enumerate}
\item 证明：\(A_1D\perp\) 平面 \(A_1BC\)；
\item 求直线 \(A_1B\) 和平面 \(BB_1C_1C\) 所成的角的正弦值.
\end{enumerate}

\bitmapfigure[max width=0.55\linewidth]{img/2015/zhejiang_liberal/q18_fig1.png}
\end{problem}

\begin{solution}
\begin{enumerate}
\item 设 \(E\) 为 \(BC\) 的中点. 因为 \(AB=AC\)，所以 \(AE\perp BC\). 又因为 \(A_1\) 在底面 \(ABC\) 的射影为 \(E\)，所以 \(A_1E\perp\) 平面 \(ABC\)，从而 \(A_1E\perp AE\). 在三棱柱中，\(ED\parallel AA_1\) 且 \(ED=AA_1\)，所以四边形 \(A_1AED\) 是平行四边形，故 \(A_1D\parallel AE\). 于是 \(A_1D\perp A_1E\)，且 \(A_1D\perp BC\). 由于 \(A_1E\) 与 \(BC\) 是平面 \(A_1BC\) 内相交的两条直线，所以 \(A_1D\perp\) 平面 \(A_1BC\).

\item 建立空间直角坐标系，取 \(E\) 为原点，令 \(BC\) 为 \(x\) 轴，\(AE\) 为 \(y\) 轴，垂直于底面且经过 \(E\) 的直线为 \(z\) 轴. 则
\(A(0,\sqrt2,0)\)，\(B(\sqrt2,0,0)\)，\(C(-\sqrt2,0,0)\).
设 \(A_1=(0,0,h)\). 由 \(A_1A=4\)，得 \(h^2+2=16\)，所以 \(h=\sqrt{14}\). 三棱柱的平移向量为 \(\overrightarrow{AA_1}=(0,-\sqrt2,\sqrt{14})\)，从而
\(A_1=(0,0,\sqrt{14})\)，\(B_1=(\sqrt2,-\sqrt2,\sqrt{14})\)，\(C_1=(-\sqrt2,-\sqrt2,\sqrt{14})\).
平面 \(BB_1C_1C\) 内有方向向量 \(\overrightarrow{BC}=(2\sqrt2,0,0)\) 和 \(\overrightarrow{BB_1}=(0,-\sqrt2,\sqrt{14})\)，可取其法向量 \(\bs{n}=(0,\sqrt{14},\sqrt2)\). 直线 \(A_1B\) 的方向向量为 \(\bs{v}=(\sqrt2,0,-\sqrt{14})\)，且 \(|\bs{v}|=|\bs{n}|=4\). 设所成的角为 \(\theta\)，则
\[
\sin\theta=\frac{|\bs{v}\cdot\bs{n}|}{|\bs{v}||\bs{n}|}=\frac{2\sqrt7}{16}=\frac{\sqrt7}{8}
\]
\end{enumerate}
\end{solution}

\begin{problem}
如图，已知抛物线 \(C_1:y=\frac{1}{4}x^2\)，圆 \(C_2:x^2+(y-1)^2=1\)，过点 \(P(t,0)\)（\(t>0\)）作不过原点 \(O\) 的直线 \(PA\)，\(PB\) 分别与抛物线 \(C_1\) 和圆 \(C_2\) 相切，\(A\)，\(B\) 为切点.

\begin{enumerate}
\item 求点 \(A\)，\(B\) 的坐标；
\item 求 \(\triangle PAB\) 的面积.
\end{enumerate}

注：直线与抛物线有且只有一个公共点，且与抛物线的对称轴不平行，则称该直线与抛物线相切，称该公共点为切点.

\begingroup
\leftskip=0pt\relax
\rightskip=0pt\relax
\bitmapfigure[max width=0.55\linewidth]{img/2015/zhejiang_liberal/q19_fig1.png}
\endgroup
\end{problem}

\begin{solution}
\begin{enumerate}
\item 由题注，直线 \(PA\) 不与抛物线的对称轴平行，因此可设其方程为 \(y=k(x-t)\). 因该直线不过原点，故 \(k\neq0\). 与抛物线方程联立，得
\[
x^2-4kx+4kt=0
\]
相切时判别式为零，即 \(16k^2-16kt=0\). 由 \(k\neq0\)，得 \(k=t\). 此时方程的唯一根为 \(x=2t\)，所以 \(A=(2t,t^2)\).

设 \(B=(x,y)\). 圆 \(C_2\) 的圆心为 \(D=(0,1)\)，相切条件为 \(\overrightarrow{DB}\perp\overrightarrow{PB}\)，于是
\[
(x,y-1)\cdot(x-t,y)=0
\]
又因 \(B\) 在圆上，满足 \(x^2+y^2-2y=0\). 将两式相减得 \(y=tx\). 代入圆方程，得
\[
x\left((1+t^2)x-2t\right)=0
\]
根 \(x=0\) 对应 \(B=O\)，会使直线 \(PB\) 经过原点，不合题意，故
\[
B\left(\frac{2t}{1+t^2},\frac{2t^2}{1+t^2}\right)
\]

\item 直线 \(PA\) 的方程为 \(tx-y-t^2=0\)，且
\[
|AP|=\sqrt{t^2+t^4}=t\sqrt{1+t^2}
\]
点 \(B\) 到直线 \(PA\) 的距离为
\[
d=\frac{|t\cdot\frac{2t}{1+t^2}-\frac{2t^2}{1+t^2}-t^2|}{\sqrt{1+t^2}}=\frac{t^2}{\sqrt{1+t^2}}
\]
所以
\[
S_{\triangle PAB}=\frac12|AP|d=\frac12\cdot t\sqrt{1+t^2}\cdot\frac{t^2}{\sqrt{1+t^2}}=\frac{t^3}{2}
\]
\end{enumerate}
\end{solution}

\begin{problem}
设函数 \(f(x)=x^2+ax+b\)（\(a\)，\(b\in\R\)）.

\begin{enumerate}
\item 当 \(b=\frac{a^2}{4}+1\) 时，求函数 \(f(x)\) 在 \([-1,1]\) 上的最小值 \(g(a)\) 的表达式；
\item 已知函数 \(f(x)\) 在 \([-1,1]\) 上存在零点，\(0\leqslant b-2a\leqslant 1\). 求 \(b\) 的取值范围.
\end{enumerate}
\end{problem}

\begin{solution}
\begin{enumerate}
\item 由条件
\[
f(x)=x^2+ax+\frac{a^2}{4}+1=\left(x+\frac a2\right)^2+1
\]
当 \(a\leqslant-2\) 时，顶点横坐标 \(-\frac a2\geqslant1\)，区间内离顶点最近的点为 \(x=1\)，故 \(g(a)=\left(1+\frac a2\right)^2+1=\frac{a^2}{4}+a+2\). 当 \(-2<a\leqslant2\) 时，顶点在 \([-1,1]\) 内，故 \(g(a)=1\). 当 \(a>2\) 时，顶点横坐标 \(-\frac a2<-1\)，区间内离顶点最近的点为 \(x=-1\)，故 \(g(a)=\left(-1+\frac a2\right)^2+1=\frac{a^2}{4}-a+2\). 因此
\[
g(a)=\begin{cases}
\dfrac{a^2}{4}+a+2,&a\leqslant-2,\\
1,&-2<a\leqslant2,\\
\dfrac{a^2}{4}-a+2,&a>2.
\end{cases}
\]

\item 设零点为 \(x\in[-1,1]\)，并令 \(s=b-2a\)，则 \(0\leqslant s\leqslant1\). 由 \(f(x)=0\)，得
\(a=-\frac{x^2+s}{x+2}\)，\(b=2a+s=\frac{x(s-2x)}{x+2}\).
对固定的 \(x\)，右式关于 \(s\) 为一次函数. 若 \(x\leqslant0\)，其最小值在 \(s=1\) 时取得，即
\[
\frac{x(1-2x)}{x+2}
\]
该式在 \([-1,0]\) 上递增，最小值为 \(-3\)；若 \(x\geqslant0\)，其最小值在 \(s=0\) 时取得，即 \(-\dfrac{2x^2}{x+2}\geqslant-\dfrac23\). 所以 \(b\) 的最小值为 \(-3\).

当 \(x\leqslant0\) 时，\(b\) 的最大值不超过 \(0\). 当 \(x\geqslant0\) 时，最大值在 \(s=1\) 时取得，令
\(h(x)=\frac{x(1-2x)}{x+2}\)，\(0\leqslant x\leqslant1\).
有
\[
h'(x)=\frac{2-8x-2x^2}{(x+2)^2}
\]
导数在 \(x=\sqrt5-2\) 左侧为正、右侧为负，所以最大值在该点取得，且
\[
h(\sqrt5-2)=9-4\sqrt5
\]
参数区域 \([-1,1]\times[0,1]\) 连通，\(b=\dfrac{x(s-2x)}{x+2}\) 连续，因此其中间所有值均能取得. 故 \(b\) 的取值范围为 \([-3,9-4\sqrt5]\).
\end{enumerate}
\end{solution}
